% Figure: a toggle press. A handle drives a pair of links that push a ram
% down against a workpiece. As the links approach the straight (dead-centre)
% line, the mechanical advantage rises without bound, which is why toggle
% presses, riveters, and rock crushers deliver large force at the end of the
% stroke for a modest handle push.
\begin{figure}[htbp]
\centering
\begin{tikzpicture}[>=latex, x=1.0cm, y=1.0cm]
  % fixed frame: top anchor pin O, vertical guide for the ram on the right
  \coordinate (O) at (0,3.0);      % top link anchor (fixed)
  \coordinate (K) at (1.4,1.6);    % knee pin (where handle pushes)
  \coordinate (R) at (0,0.2);      % ram pin, constrained to vertical line x=0
  % frame outline (top bracket)
  \draw[engblack, thick] (-0.8,3.0) -- (0.8,3.0);
  \foreach \x in {-0.7,-0.5,...,0.7} { \draw[enggray, thin] (\x,3.0) -- (\x-0.14,3.18); }
  % vertical guide slot for the ram
  \draw[enggray, thick] (-0.35,-0.2) -- (-0.35,1.4);
  \draw[enggray, thick] (0.35,-0.2) -- (0.35,1.4);
  % two links: O--K and K--R
  \draw[engblue, very thick] (O) -- (K);
  \draw[engblue, very thick] (K) -- (R);
  % knee angle theta each link makes with the line of travel (vertical)
  \draw[enggray, thin] ($(K)+(-0.5,0.55)$) arc (250:290:0.5);
  \node[enggray, font=\scriptsize] at ($(K)+(-0.05,0.85)$) {$\theta$};
  % handle push P at the knee, horizontal inward
  \draw[->, very thick, engorange] ($(K)+(1.5,0)$) -- (K);
  \node[engorange, font=\scriptsize, anchor=west] at ($(K)+(0.6,0.25)$) {$P$};
  % ram and clamp force Q downward onto workpiece
  \draw[engblue, very thick, fill=engblue!12] (-0.32,0.2) rectangle (0.32,-0.35);
  \draw[->, very thick, engred] (0,-0.35) -- ++(0,-0.7);
  \node[engred, font=\scriptsize, anchor=west] at (0.05,-0.75) {$Q$};
  % workpiece
  \draw[engblack, thick, fill=enggray!25] (-0.6,-1.3) rectangle (0.6,-1.05);
  \node[font=\scriptsize, anchor=north] at (0,-1.3) {workpiece};
  % joints
  \fill[engblack] (O) circle (1.8pt); \node[font=\scriptsize, anchor=south east] at (O) {$O$};
  \fill[engblack] (K) circle (1.8pt); \node[font=\scriptsize, anchor=south west] at (K) {knee};
  \fill[engblack] (R) circle (1.8pt); \node[font=\scriptsize, anchor=east] at ($(R)+(-0.4,0.1)$) {ram};
  % dead-centre reference: dotted vertical through O and R
  \draw[englime!55!black, dotted] (O) -- (R);
  \node[englime!55!black, font=\scriptsize, anchor=west] at (0.05,2.3) {dead-centre line};
\end{tikzpicture}
\caption[Toggle press]{A toggle press. A handle push $P$ at the knee drives two
links, $OK$ anchored to the frame and $KR$ driving the ram, toward the straight
dead-centre line joining the fixed anchor $O$ to the ram $R$. Let $\theta$ be
the angle each link makes with the line of travel. A small handle motion
advances the ram through the link geometry, and the virtual-work balance gives a
clamp force $Q$ that scales as $P\cot\theta$: as the links straighten and
$\theta\to 0$ the advantage runs toward infinity, so the press delivers its
greatest force at the very end of the closing stroke. The same geometry that
multiplies force starves travel, so the ram moves least where it pushes hardest,
the force-distance trade every machine keeps. Riveters, rock crushers, injection
clamps, and knee-action switches all live on this rising advantage near
dead-centre, and the slight over-centre setting that many designs adopt lets the
clamp force itself hold the mechanism closed.}
\label{fig:vol03ch02:toggle-press}
\end{figure}
